\documentclass[10pt,a4paper]{amsart}
\usepackage{amsmath,amssymb,amsthm,graphicx}
\RequirePackage{amssymb,amsmath,amsthm,bm}
\RequirePackage{thmtools}
\RequirePackage[capitalize]{cleveref}
\providecommand{\tightlist}{%
  \setlength{\itemsep}{0pt}\setlength{\parskip}{0pt}}

%%%%% varphi instead of phi
\renewcommand{\phi}{\varphi}

\mathchardef\mhyph="2D 
\providecommand{\A}{\mathcal{A}}
\providecommand{\C}{\mathcal{C}}
\providecommand{\L}{\mathcal{L}}
\providecommand{\B}{\mathcal{B}}
\providecommand{\M}{\mathcal{M}}
\renewcommand{\S}{\mathcal{S}}
\providecommand{\ol}[1]{\overline{#1}}
\providecommand{\mbf}[1]{\mathbf{#1}}
\providecommand{\mf}[1]{\mathfrak{#1}}
\providecommand{\mrm}[1]{\mathrm{#1}}
\providecommand{\mc}[1]{\mathcal{#1}}
\providecommand{\ra}{\rightarrow}
\providecommand{\Ra}{\Rightarrow}
\providecommand{\La}{\Leftarrow}
\providecommand{\lr}{\leftrightarrow}
\providecommand{\hf}{\mathbb{HF}}
\providecommand{\spec}{DgSp}
\providecommand{\bSigma}{\bm\Sigma}
\providecommand{\bDelta}{\bm\Delta}
\providecommand{\bPi}{\bm\Pi}
\providecommand{\Pinf}[1]{\Pi^{\mathrm{in}}_{#1}}
\providecommand{\Sinf}[1]{\Sigma^{\mathrm{in}}_{#1}}
\providecommand{\dSinf}[1]{d\mhyph\Sigma^{\mathrm{in}}_{#1}}
\providecommand{\Dinf}[1]{\Delta^{\mathrm{in}}_{#1}}
\providecommand{\bp}{\bar p}
\providecommand{\bq}{\bar q}
\providecommand{\ba}{\bar a}
\providecommand{\bb}{\bar b}
\providecommand{\bc}{\bar c}
\providecommand{\bd}{\bar d}
\providecommand{\bx}{\bar x}
\providecommand{\by}{\bar y}
\providecommand{\bz}{\bar z}
\providecommand{\LO}{\mathrm{LO}}
\declaretheorem[name={Theorem},numberwithin=section]{theorem}
\declaretheorem[name={Proposition}, sibling=theorem]{proposition}
\declaretheorem[name={Lemma}, sibling=theorem]{lemma}
\declaretheorem[name={Fact}, sibling=theorem]{fact}
\declaretheorem[name={Claim}, numberwithin=theorem]{claim}
\declaretheorem[name={Definition},numberwithin=section,sibling=theorem, style=definition]{definition}
\declaretheorem[name={Remark}, numbered=no,style=definition]{remark}
\declaretheorem[name={Corollary}, sibling=theorem]{corollary}
\declaretheorem[name={Question},sibling=theorem]{question}
\newcommand{\T}{{\mathcal{T}}}
\newcommand{\SR}{SR}
\newcommand{\pSR}{pSR}
\newcommand{\Q}{\mathbb{Q}}
\newcommand{\Z}{\mathbb{Z}}
\newcommand{\K}{\mathcal{K}}
\newcommand{\N}{\mathcal{N}}
\renewcommand{\L}{\mathcal{L}}
\providecommand{\Ginf}[1]{\Gamma^{\mathrm{in}}_{#1}}
\title{Formulas over a structure translate to formulas over its labelled tree at the same quantifier complexity, and back}
\author{David Gonzalez, Dino Rossegger}
\date{arxiv:2305.07126v1; CC BY 4.0}
\begin{document}
\maketitle
\noindent\textit{Editorial note.} The counted statement is the author environment at source0.tex lines 234-249 and its proof environment opens at line 250, the line immediately after; the pairing is by adjacency and the retained proof range 250-286 is the whole environment.\par\medskip
\noindent\textit{Editorial note.} The retained preamble is the author style file dino.sty (package lines, the theorem declaration support, the notation and tuple macros and the theorem environments) followed by the ordinal and rank macros of the manuscript itself, reproduced verbatim; this is the same recipe as the previously published card of this article, so no package outside it is requested and extra\_packages stays empty.\par\medskip
\noindent\textit{Editorial note.} Metadata provenance: source\_id and source\_version are the values recorded in this source licence evidence file (sources/scott-linear-orders-2305.07126v1/LICENSE\_EVIDENCE.json: arxiv:2305.07126v1 / 2305.07126v1), which the card points at directly; the title and author list are transcribed from the authors\_index and title keys of the source record of the pinned copy. The card carries no counter setting and imposes no original numbering, since this article has no label-number dump and its numbering cannot be reconstructed from the retained ranges alone.\par\medskip
\noindent\textit{Editorial note.} Two item labels are defined inside the retained statement (eq:pushforward and eq:pullback) and are the two directions the retained proof refers to; they need no external entry, and the retained theorem and proof use no label from outside the retained range.\par\medskip
\noindent\textit{Editorial scope.} One original result of the pinned author manuscript, counted once: the two-way translation between formulas over a structure and formulas over its labelled tree at equal quantifier rank (author label lemma:translation).\par\medskip
\noindent\textit{Allowed standard prerequisites (editorial).}\begin{itemize}
\item Labelled tree of a structure. Every structure A has a labelled tree T\_A of finite tuples ordered by extension, and for every tuple a in A there is a node sigma of T\_A above it.. It is the object the two translations go between.
\item Infinitary formulas of bounded rank. P\_alpha and the classes Gamma\textasciicircum{}in\_alpha are the infinitary formulas of quantifier rank below alpha built with the corresponding operations.. They are the classes the two directions preserve.
\item Turing computable embedding. A Turing computable embedding from a class of structures with parameters into labelled trees with parameters, extending the parameter-free Friedman-Stanley embedding.. It carries the pull-back direction of the proof.
\end{itemize}\medskip
\begin{lemma}\label{lemma:translation} 
   Let $\A$ be a structure in vocabulary $\tau$.
   \begin{enumerate}
      \tightlist
      \item\label{eq:pushforward} Given a $\Pinf{\alpha}$ $\tau$-formula $\phi$, there is a
         $\Pinf{\alpha}$ formula $T_\phi$ in the vocabulary of labeled trees
         such that for any $\ba\in \A$ and every $\sigma\in \T_\A$ with
         $\sigma\succeq \ba$ 
         \[\A\models\phi(\ba)\iff \T_\A\models T_\phi(\sigma).\]
      \item\label{eq:pullback} Given a $\Pinf{\alpha}$ formula $\psi$ in the
         vocabulary of labeled trees, there is a $\Pinf{\alpha}$ $\tau$-formula
         $\psi^*$ such that for any tuple $\bar{\sigma}=(\ba_1,\dots,\ba_n)\in
         \T_\A$ \[\T_\A\models \psi(\bar\sigma) \iff
         \A\models\psi^*(\ba,\dots,\ba_n).\]
      \end{enumerate}
\end{lemma} 

\begin{proof}
   We first show \cref{eq:pushforward} starting with quantifier free formulas.
   Let $\phi$ be quantifier free with free variables $x_1,\dots x_n$ in
   disjunctive normal form. For each disjunct $\phi_i$ let $n_i$ be the index of
   the largest occuring variable and let $T_{\phi_i}$ be the formula expressing
   that the height of $x$ in the tree is at least $n_i$ and that
   $\phi_i\subseteq l(x)$. Then $T_\phi=\bigvee T_{\phi_i}$. The construction of
   $\T_\A$ guarantees that for $\sigma\succeq\ba\in \A$, $\T_\A\models
   T_\phi(\sigma)$ if and only if $\A\models \phi(\ba)$. Now, assume that we
   have a translation for $\phi(\bx,y)$ into $T_\phi(x)$. If $\psi=\exists y
   \phi(\bx,y)$, then let \[T_\psi= (\exists y\succeq x)\, T_\phi(y).\] We have
   that $\A\models \exists y\,\psi(\ba,y)$, if and only if there is $b\in \A$
   such that $\A\models\psi(\ba,b)$ if and only if there is $\sigma\succeq \ba
   \in \T_\A$ $\T_\A\models T_\psi(\sigma)$ if and only if $\T_\A\models
   T_\psi$. This proves the successor case. The limit case is immediate.

   We now prove \cref{eq:pullback} by defining a family of Turing-computable
   embeddings from the class of structures containing $\A$ with parameters to
   labeled trees with parameters. Each of these embeddings will in some sense
   extend the Friedman-Stanley embedding that acts without parameters. Fix a
   tuple $\bar\sigma=(\ba_1,\cdots,\ba_n)\in \T_\A$. Consider any structure $\B$
   in the class of $\A$ with parameters $\bp$ such that
   $|\bp|=|\ba_1\dots\ba_n|$. We now construct the embedding
   $\Psi_{\bar\sigma}$. Non-uniformly fix a partition
   $\bp=\bp_1,\dots\bp_n$ with $|\bp_i|=|\ba_i|$ and construct the tree of
   tuples of $\B$ as normal. To decide what elements become parameters in
   $\T_\B$, look for the lexicographical least elements $\tau_i$ coding the
   $\bp_i$ such that the finite substructure of $\T_\B$ described by $\tau$ is
   isomorphic to the one described by $\sigma$ in $\T_\A$. It is not difficult to check that this is a Turing computable
   embedding. Now, for any formula $\psi$ in the language of labeled trees, any
   $\B$ in the class of $\A$ and any $\bar\sigma=(\bb_0\dots\bb_n)\in \T_\B$ by
   the pull-back theorem for $\Psi_{\bar\sigma}$ there is $\psi^*$ of the same
   complexity as $\psi$ such that $\B\models \psi^*(\bb_0\dots\bb_n)$ if and
   only if $\T_\B\models \psi(\bar\sigma)$. As the image of $\Psi_{\bar\sigma}$
   contains the parameterized tree $(\T_\A,\bar\sigma)$, this proves
   \cref{eq:pullback}.
\end{proof}

\end{document}
